\documentclass[11pt]{article}

\usepackage[margin=1in]{geometry}
\usepackage{amsmath,amssymb}
\usepackage{bm}
\usepackage[hidelinks]{hyperref}

\newcommand{\kv}{\mathbf{k}}
\newcommand{\qv}{\mathbf{q}}
\newcommand{\term}[1]{\item[#1.]}

\title{\texttt{ironmc} / \texttt{rmc-ab-initio} project glossary}
\author{Samuel Recker}
\date{\today}

\begin{document}
\maketitle

\noindent
Working vocabulary of the ab-initio DiagMC project. Terms are grouped into
\textbf{physics} (Section~\ref{sec:physics}) and \textbf{software engineering}
(Section~\ref{sec:swe}); overloaded words that mean different things in different
contexts are disambiguated in Section~\ref{sec:overloaded}. The methodology note
\texttt{latex/ab-initio-diagmc/ab-initio-diagmc.tex} is the authoritative physics
reference; this file only defines the words.

\section{Physics}
\label{sec:physics}

\subsection{Model and observables}

\begin{description}
\term{Polaron} An electron dressed by the lattice distortion (phonon cloud) it drags
  along. The project computes the \emph{electron} polaron of LiF: one carrier at the
  conduction-band minimum plus its phonon cloud.

\term{Formation energy $E_{0}$} The polaron ground-state energy relative to the bare
  band edge, $E_{0}=E_{\mathrm{polaron}}-\varepsilon_{\mathrm{CBM}}$. Negative;
  the headline observable of every production run. LiF references:
  $-0.2479$~eV on the $20^{3}$ grid, $-0.408$~eV extrapolated (FEP-DMC).

\term{Fr\"ohlich model} The textbook continuum polaron: parabolic electron band, one
  dispersionless LO phonon, vertex $|g(\qv)|^{2}\propto\alpha/q^{2}$. Serves two roles
  here: the analytic limit LiF is close to (so an analytic cross-check exists), and the
  engine-validation model (see \emph{stage}, Section~\ref{sec:swe-validation}).

\term{Coupling constant $\alpha$} The dimensionless Fr\"ohlich electron--phonon
  coupling. $\alpha_{\mathrm{LiF}}\approx4.76$ (intermediate-to-strong coupling),
  computed from $m^{*}$, $\omega_{\mathrm{LO}}$, $\varepsilon_{\infty}$,
  $\varepsilon_{0}$. Weak-coupling target: $E^{(2)}\approx-\alpha\hbar\omega_{\mathrm{LO}}$.

\term{LO / TO phonons, LO--TO splitting} Longitudinal / transverse optical branches. In
  a polar crystal the long-range dipole field stiffens the LO branch at $\Gamma$
  (LiF: LO $18.45$~THz vs TO $8.12$~THz). The splitting is carried entirely by the
  long-range part of the dynamical matrix, which EPW subtracts and the code must restore.

\term{Imaginary time $\tau$} The evolution variable of the zero-temperature
  imaginary-time formulation. With $\hbar=1$ and energies in eV, $\tau$ has units
  eV$^{-1}$. Diagrams live on a $\tau$ interval $[0,\tau_{\max}]$.

\term{Green's function $G(\tau)$} The imaginary-time electron propagator whose
  large-$\tau$ decay $G(\tau)\sim Z e^{-(E_{0}-\mu)\tau}$ encodes the ground-state
  energy. What the decay estimator fits; \emph{not} the same object as the self-energy.

\term{Chemical potential $\mu$} A reweighting shift $e^{+\mu\tau}$ applied to the
  electron propagator to keep diagram weights bounded. In the decay estimator it must
  bracket $E_{0}$ from below (P9 of the methodology note); in the fixed-$\tau$
  thermodynamic ensemble it drops out entirely.

\term{Units (physical vs natural)} The ab-initio model runs in physical units:
  energies eV, time eV$^{-1}$, momenta in fractional (crystal) coordinates on a
  $\Gamma$-centered grid. Fr\"ohlich toy runs use polaron natural units
  $\hbar=m=\omega_{\mathrm{LO}}=1$. Ab-initio inputs arrive in Rydberg
  ($\approx13.6057$~eV) or THz and are converted on load.

\term{Effective mass $m^{*}$, CBM} Conduction-band minimum and its curvature mass.
  LiF's CBM is single and non-degenerate at $\Gamma$ (why LiF is the first material);
  $m^{*}$ enters $\alpha$ and the $q$-unit bridge of the reduced Fr\"ohlich comparison.

\term{Second-order energy $E^{(2)}$} The lowest-order (one-phonon) perturbative
  polaron energy, $E^{(2)}(\kv)=-\frac{1}{N_{q}}\sum_{\qv\nu}
  |g(\kv,\qv,\nu)|^{2}/(\varepsilon_{\kv+\qv}-\varepsilon_{\kv}+\omega_{\qv\nu})$.
  Cheap to evaluate exactly; used as the vertex-scale check and as a (biased) scanning
  proxy for compression error --- it overestimates the in-sampler systematic by
  $3$--$10\times$, so it never decides a budget (see \emph{gate}).

\term{Lang--Firsov / dispersionless model} Exactly solvable lattice polaron
  ($\varepsilon(\kv)=\mathrm{const}$), $E_{0}=-\bar g^{2}/\omega$ at all couplings.
  Used as an analytic validation target for the thermodynamic estimator.

\term{Discretized Hamiltonian $\hat H_{n}$} The polaron Hamiltonian restricted to a
  uniform $n^{3}$ Brillouin-zone grid shared by electron and phonon momenta, closed
  under momentum addition, with per-arc measure $\frac{1}{N_{q}}\sum_{\qv\neq\Gamma}$,
  $N_{q}=n^{3}$. \emph{This} is what the sampler solves --- numerically exactly. All
  physics statements are statements about the sequence $E_{0}(n)$.

\term{Thermodynamic limit / $1/n$ extrapolation} Because the discrete BZ sum
  underweights the $1/q^{2}$ polar region near $\Gamma$, $E_{0}(n)$ converges from
  above with leading error $\propto1/n$. The continuum answer is the intercept of a
  linear fit of $E_{0}(n)$ against $1/n$ over the grid ladder $n=20,32,40,50$
  (same methodology as FEP-DMC).
\end{description}

\subsection{Diagrammatics and sampling}

\begin{description}
\term{DiagMC (diagrammatic Monte Carlo)} Stochastic sampling of the Feynman-diagram
  expansion of an observable --- here the imaginary-time worldline expansion of the
  polaron partition function / Green's function --- by a Metropolis Markov chain over
  diagram topologies, momenta, and times.

\term{Diagram, worldline, backbone} One Monte-Carlo configuration: an electron line
  (backbone) on $[0,\tau]$, cut into \emph{segments} by interaction vertices, with
  phonon \emph{arcs} connecting vertex pairs. Stored as a dense $\tau$-sorted array
  (GPU) or slotmap (CPU).

\term{Segment} A piece of electron line between adjacent vertices, carrying momentum
  $\kv_{s}$ and weight $e^{-\varepsilon(\kv_{s})\Delta\tau_{s}}$.

\term{Arc (phonon arc)} A phonon propagator $D(\omega_{\qv\nu},\Delta\tau)$ joining
  two vertices; carries branch $\nu$ and momentum $\qv$, which is subtracted from the
  backbone momentum underneath it.

\term{Vertex (diagrammatic)} An electron--phonon interaction event at a time
  $\tau_{i}$, weighted by the matrix element $g(\kv,\qv,\nu)$. Not to be confused with
  ``vertex'' as shorthand for the matrix element itself (common in this repo).

\term{Diagram order $n$, $\langle n\rangle$ (\texttt{mean\_order})} Number of phonon
  arcs in the current diagram. The equilibrium mean order tracks coupling strength
  (LiF at $20^{3}$: $\langle n\rangle\approx16$); a collapsing or runaway order is the
  standard symptom of a mis-set $\mu$, a mis-scaled vertex, or under-equilibration.

\term{Update} One Metropolis move proposal (insert/remove arc, shift times, swap,
  \dots). The engine is built from an update set run by a \texttt{Runner} through a
  \texttt{MetropolisKernel}.

\term{Chain} One independent Markov chain with its own deterministic RNG stream.
  Production runs use many chains ($12$ per CPU node config, $8192$ per GPU) for
  honest, embarrassingly parallel statistics.

\term{Warmup / equilibration transient} Steps discarded before measuring, which must
  exceed the diagram-order equilibration time ($\sim10^{5}$ steps for LiF at
  $20^{3}$) \emph{per chain} --- many short chains do \emph{not} dilute the transient
  the way one long chain does (the GPU warmup lesson).

\term{Sign, average sign $\langle s\rangle$} The complex vertex makes diagram weights
  non-positive-definite; each sample carries $\operatorname{sign}(\operatorname{Re}\,w)$
  and observables are divided by $\langle s\rangle$. For LiF(e) the sign stays
  benign ($\approx0.99$); for the Fr\"ohlich model it is $\equiv1$.

\term{$|g|^{2}$-weighted proposal} The native phonon-momentum proposal distribution:
  new arc momenta are drawn $\propto|g(\qv)|^{2}$ (tabulated), so the polar $1/q^{2}$
  peak is proposed where the weight lives instead of rejected.
\end{description}

\subsection{Estimators and statistics}

\begin{description}
\term{Decay estimator} Reads $E_{0}$ from the large-$\tau$ slope of the measured
  $G(\tau)$, $E_{0}=\mu+\mathrm{slope}$, fitted beyond \texttt{tau\_fit}. Requires
  bracketing $\mu$ and a fit window clear of excited-state contamination
  ($\tau_{\mathrm{fit}}\gg1/\omega_{\mathrm{LO}}$); acceptance criterion is
  $\mu$-independence of the result. Superseded for production by the thermodynamic
  estimator; kept as a cross-check.

\term{Thermodynamic (fixed-$\tau$) estimator} The FEP-DMC-style operator estimator:
  the sampler runs in a fixed-$\tau$ ensemble (external time pinned at
  $\tau_{\max}=\beta$, \texttt{fixed\_tau: true}) and averages the local energy
  $E_{\mathrm{loc}}=\tau^{-1}[\sum_{\mathrm{seg}}\varepsilon\,\Delta\tau
  +\sum_{\mathrm{arcs}}\omega\,\Delta\tau-2N_{\mathrm{arcs}}]$ over the whole
  worldline --- no fit, no $\mu$ bracketing, lower variance. Reported as
  \texttt{e0\_thermo}. Systematic error controlled by $\beta$
  (excited-state contamination $\sim e^{-\Delta\beta}$), checked by $\tau$-doubling
  (guard~A).

\term{Exp-reweighted estimator, \texttt{energy\_estimate}} The older stage-1/stage-2
  estimator $\langle-\exp((E_{\mathrm{est}}-\mu)\tau)\rangle$ whose self-consistent
  update diverges for $\alpha\gtrsim1$ (fixed-point slope $\approx-2.2$); the stage
  configs freeze \texttt{energy\_estimate} and it is bisected by hand to the fixed
  point $E\approx E_{\mathrm{est}}$.

\term{Local energy} The per-worldline value of $\langle\hat H\rangle$ decomposed into
  band, phonon, and interaction accumulators (segment $\varepsilon\Delta\tau$ sums,
  arc $\omega\Delta\tau$ sums, vertex count); its sign-weighted time average is the
  thermodynamic estimator.

\term{Thermal phonon propagator}
  $D(\omega,\tau)=[e^{-\omega\tau}+e^{-\omega(\beta-\tau)}]/(1-e^{-\omega\beta})$:
  spontaneous emission $+$ thermal wrap-around, Bose-enhanced. The
  \texttt{zeroTMC} limit hard-codes $\beta\to\infty$, i.e.\ $D=e^{-\omega\tau}$;
  LiF(e) production runs are effectively $T=0$. Modes below a
  $\sim1$~meV frequency cutoff are dropped to avoid the $\omega\to0$ divergences.

\term{Jackknife} Leave-one-out resampling over measurement batches for error bars,
  computed as total$-$batch (never re-aggregating per replicate --- the $O(N^{2})$
  trap). With independent chains the natural batching is one batch per chain.

\term{Batch (\texttt{n\_batches})} A contiguous block of measurements forming one
  jackknife unit. Under deadline truncation the batch pooling still guarantees
  $\geq$ one block per chain.

\term{\texttt{steps\_done}} The actually executed step count reported by a run; all
  quoted steps/s throughput numbers use it, so deadline-truncated runs stay honest.
\end{description}

\subsection{Ab-initio pipeline (DFT $\to$ DFPT $\to$ Wannier $\to$ vertex)}

\begin{description}
\term{Quantum ESPRESSO (QE)} The plane-wave DFT suite: \texttt{pw.x} (SCF/NSCF
  ground state and band structure), \texttt{ph.x} (DFPT phonons). One QE version must
  be used for the whole \texttt{scf/ph/nscf/epw} chain.

\term{DFPT} Density-functional perturbation theory; supplies phonon frequencies and
  eigenvectors, Born effective charges $\mathbf{Z}^{*}$, and the electronic dielectric
  tensor $\bm\varepsilon_{\infty}$.

\term{EPW} The QE companion code that Wannier-interpolates the electron--phonon
  coupling; produces the Wannier-representation tensor \texttt{LiF.epmatwp} and (with
  \texttt{lpolar=.true.}) subtracts the polar long-range part first.

\term{Wannier functions / interpolation} Localized real-space orbitals obtained from
  the Bloch states; matrix elements stored on real-space pairs $(\mathbf{R}_{e},
  \mathbf{R}_{p})$ interpolate to any $(\kv,\qv)$ by Fourier transform. The electron
  dispersion $\varepsilon(\kv)$ used by the sampler is the Wannier-interpolated
  conduction band (not a parabola).

\term{Disentanglement windows} Wannierization energy windows
  (\texttt{dis\_win\_max}, \texttt{dis\_froz\_max}) selecting/freezing bands. Must be
  re-tuned when pseudopotentials shift the band edges, or the frozen CBM slips out.

\term{Pseudopotentials: ultrasoft vs norm-conserving (ONCV)} EPW's e--ph matrix
  elements require \emph{norm-conserving} pseudopotentials; the initial ultrasoft
  (GBRV) deck silently produced an all-NaN tensor (pitfall P1). The fix (Pseudo Dojo
  ONCV) raises \texttt{ecutwfc} $50\to96$~Ry.

\term{\texttt{.epmatwp}} EPW's Wannier-representation e--ph tensor file
  $g(\mathbf{R}_{e},\mathbf{R}_{p})$ --- the primary input the Rust reader consumes.

\term{\texttt{.epb}} EPW's coarse Bloch-representation dump: the \emph{full},
  pre-subtraction coupling at the coarse $q$-points, plus eigenvalues, dynamical
  matrices, $\mathbf{Z}^{*}$, $\bm\varepsilon_{\infty}$. One Fortran unformatted
  record per MPI pool; the ground truth for the round-trip validation.

\term{\texttt{dynq}} The coarse DFPT dynamical matrices; ground truth for the phonon
  (P7) round-trip.

\term{Vertex $g_{mn\nu}(\kv,\qv)$} The electron--phonon matrix element: scattering of
  an electron from band $n$ at $\kv$ to band $m$ at $\kv+\qv$ emitting/absorbing a
  phonon of branch $\nu$, momentum $\qv$. Complex; the sampler consumes
  gauge-invariant $|g|^{2}$ combinations.

\term{Short/long-range split $g=g^{\mathcal S}+g^{\mathcal L}$} EPW subtracts the
  analytic polar (Fr\"ohlich-like, $1/q$) part $g^{\mathcal L}$ before Wannierization
  so the stored tensor $g^{\mathcal S}$ is short-ranged; the code re-adds
  $g^{\mathcal L}$ (Verdi--Giustino form, Ewald-summed) exactly, never compressed.
  The split lives in the Cartesian displacement basis and carries an electron
  band-rotation factor (pitfall P6).

\term{Polar restoration of the dynamical matrix} The same subtraction applies to the
  dynamical matrix: the rigid-ion dipole--dipole term (Gonze form, EPW
  \texttt{rgd\_blk}) must be added back before diagonalization or the LO--TO
  splitting is lost and the zero-point factor mis-scales $|g|^{2}$ by $\sim2$
  (pitfall P7).

\term{Zero-point factor} The $(\hbar/2M_{\kappa}\omega_{\qv\nu})^{1/2}$ phonon
  quantization factor in the branch projection of the coupling. Omitting it (pitfall
  P8) inflated every $|g|$ by $\sim7\times$.

\term{Acoustic sum rule} The $\qv=0$ self-term subtraction on the diagonal blocks of
  the long-range dynamical matrix that keeps the acoustic branches gapless at
  $\Gamma$.

\term{Born effective charges $\mathbf{Z}^{*}_{\kappa}$, $\bm\varepsilon_{\infty}$,
  $\varepsilon_{0}$} Per-atom dynamical charge tensors and the electronic /static
  dielectric constants; the ingredients of both long-range terms. Consistency check:
  Lyddane--Sachs--Teller,
  $\varepsilon_{0}/\varepsilon_{\infty}=(\omega_{\mathrm{LO}}/\omega_{\mathrm{TO}})^{2}$.

\term{Gauge / band rotation $U$} Bloch states carry an arbitrary per-$\kv$ phase (and
  mixing, where degenerate); the factor $[U_{\kv+\qv}U^{\dagger}_{\kv}]_{mn}$ rotates
  the polar term into the smooth Wannier gauge. Per-band comparisons are only
  meaningful in gauge-invariant combinations.

\term{$W$ star} The star of the $W$ point of the fcc BZ, contained in the $8^{3}$
  coarse grid, where the two lowest conduction bands are exactly degenerate --- making
  per-band round-trip checks gauge-ambiguous there (those $q$ are skipped; the
  sampler is unaffected).

\term{Reduced model} The single-band, $\Gamma$-slice scalar reduction
  $|g(\Gamma,\qv,\nu)|^{2}$ cached to \texttt{reduced-model.json} for the early
  stage-2 runs; superseded by the full-$\kv$ tabulated model but still built by the
  \texttt{model} subcommand.
\end{description}

\subsection{Compression and tabulation}

\begin{description}
\term{SVD compression / rank budget} The short-range Wannier tensor is compressed by
  per-band-pair truncated SVD; the \emph{global rank budget} $r$ keeps the largest
  singular values across all band pairs (production: $800$ of $1600$). A vertex
  evaluation is then a rank-$r$ contraction
  $g=\sum_{g}s_{g}u_{g}(\kv)v_{g}(\qv,\nu)+g^{\mathcal L}$.

\term{Factor tables $u,v$} The SVD factors tabulated on the momentum grid,
  stored $\kv$-major / $(\qv,\nu)$-major so one evaluation streams contiguous memory
  (the $11$--$19\times$ cache-layout win); f32 on the GPU tables, with all weight
  arithmetic f64.

\term{$\Delta_{\mathrm{rank}}$} The compression systematic,
  $E_{0}(\text{budget})-E_{0}(\text{uncompressed})$, measured by two
  otherwise-identical DMC runs \emph{in the sampler itself} and quoted on every
  production number. Current certification: $\Delta_{\mathrm{rank}}(800)=1.0\pm0.8$~meV.
  Never decided by element-wise reconstruction error or the $E^{(2)}$ proxy, and
  ``two nearby budgets agree'' is not evidence (convergence in the budget is
  non-monotonic).

\term{Uncompressed reference} The rank-$1600$ (full-rank) table used as the
  comparison baseline for $\Delta_{\mathrm{rank}}$; built on demand by
  \texttt{ironmc tables reference}.

\term{Table (\texttt{.rmcvtab})} A self-contained ladder artifact
  (\texttt{table-n\{grid\}-r\{rank\}-f32.rmcvtab}): the complete discretized
  Hamiltonian $\hat H_{n}$ --- Wannier CBM dispersion, all six polar-restored phonon
  branches, exact long-range vertex, SVD factor tables, gauge overlaps. DMC runs
  consume a table; building it is a separate, cheaper step
  (\texttt{ironmc tables build/ladder}).
\end{description}

\section{Software engineering}
\label{sec:swe}

\subsection{Validation vocabulary}
\label{sec:swe-validation}

\begin{description}
\term{Validation stages} The general pattern: correctness is certified in ordered
  stages of increasing physical content (or increasing hardware realism, for the
  GPU port), each stage mechanical to check, so a failure localizes to the first
  stage that breaks. The main sequence: engine on analytic Fr\"ohlich $\to$ vertex
  scale $\to$ compression convergence $\to$ material comparison. The GPU sequence:
  executable spec $\to$ CPU-runtime bit parity $\to$ hardware accept-count parity
  $\to$ physics anchor.

\term{Stage} One level of a validation sequence. Concretely, in run configs and the
  journal:
  \textbf{stage 1} $=$ the engine validated on the analytic Fr\"ohlich model at
  $\alpha_{\mathrm{LiF}}=4.76$ against Mishchenko \emph{et al.}\ (2000)
  (\texttt{configs/stage1-frohlich-lif.json}; sanity anchors $E_{0}(\alpha{=}1)=-1.013$,
  $E_{0}(\alpha{=}2)=-2.08$);
  \textbf{stage 2} $=$ the same engine driven by the \emph{real} LiF tensor, validated
  against literature ab-initio numbers
  (\texttt{configs/stage2-tensor-lif.json}; Sio--Giustino 2019, and concretely the
  FEP-DMC $20^{3}$ anchor). ``Stage-2-grade number'' means: good enough to survive the
  stage-2 checkpoint, i.e.\ quotable physics.

\term{Stage-2 checkpoint} The concrete pass/fail criterion of stage 2:
  $|E_{0}(20^{3})-(-0.2479~\mathrm{eV})|<30$~meV against FEP-DMC
  (\texttt{ironmc ladder checkpoint}). Current status: CPU
  $-0.24996\pm0.00026$~eV passes at $\sim2$~meV; the first GPU point sat $1.7$~meV
  shallow from under-warmup and is not yet stage-2-grade.

\term{Grid ladder} The D1/D7 production campaign: the sequence of grids
  $n=20,32,40,50$ whose $E_{0}(n)$ feed the $1/n\to0$ extrapolation. Each grid is one
  \emph{point} of the ladder (``the $20^{3}$ point''); a repo-wide rename retired the
  earlier word for this, see \emph{rung} in Section~\ref{sec:overloaded}.

\term{Ladder run} One production DMC run at a single grid point of the ladder:
  \texttt{ironmc ladder run <n> [--rank 800] [--tau120] [--node]} loads the prebuilt
  \texttt{.rmcvtab} table for $n^{3}$, runs the fixed-$\tau$ thermodynamic-estimator
  config, and returns a structured record
  \{\texttt{grid, rank, tau, e0, stderr, sign, order, steps\_done}\}. Records pipe
  into \texttt{ironmc ladder fit} for the $1/n$ extrapolation.

\term{Gate} An \emph{encoded}, numeric pass/fail acceptance check that decides a
  promotion (of a rank budget, a dataset, a backend) --- the antidote to eyeballing.
  The decided gates are cargo aliases / \texttt{ironmc gates} commands:
  \texttt{gate-e2} (the $E^{(2)}$ compression gate at $8^{3}$:
  $|E^{(2)}(800)-\mathrm{ref}|<3$~meV, $|E^{(2)}(300)-\mathrm{ref}|<12$~meV, short
  fraction $<35\%$) and \texttt{gate-dmc} (the in-sampler rank gate, budget vs the
  uncompressed $1600$ reference; scout-sized by default, production via
  \texttt{IRONMC\_DMC\_STEPS=100000000}). Thresholds are calibrated once, recorded in
  \texttt{journal.md}, then enforced mechanically.

\term{Guard} A gate run \emph{on the production configuration itself} before the
  expensive sweep: \textbf{guard A} (\texttt{ladder guard-tau}) doubles
  $\tau_{\max}$ $60\to120$ at $20^{3}$ and requires the $E_{0}$ shift $<2\sigma$
  (finite-$\beta$ systematic under control); \textbf{guard B}
  (\texttt{ladder guard-rank}) re-measures $\Delta_{\mathrm{rank}}(800)$ on the
  production grid with tolerance $\max(2\sigma,3~\mathrm{meV})$. The remaining points
  (\texttt{ladder sweep}) run only after both guards pass.

\term{Checkpoint} A comparison of a computed number against an external reference
  with a stated tolerance band (e.g.\ the stage-2 $\pm30$~meV band). Unlike a gate it
  references literature, not an internal baseline.

\term{Smoke test} A deliberately tiny end-to-end run whose only job is to prove the
  whole path executes and produces sane values before compute is spent --- e.g.\
  \texttt{ladder-lif-gpu-smoke.json} ($256$ chains $\times$ $20$k steps vs the
  production $8192\times160$k $+$ $300$k warmup), the first $8^{3}$ DMC smoke, or
  the reader's physics smoke tests. Institutional memory of pitfall P2: a run can
  finish cleanly, at the right file size, and be $100\%$ NaN --- so smoke tests check
  \emph{values}, never existence or size.

\term{Pin (pinned value)} A hard-coded expected number in a test (band energies,
  phonon frequencies, tensor dimensions, $\mathbf{Z}^{*}$ components) asserted
  against freshly-read data; ``reader pins updated'' means the dataset changed and
  the expected values were re-recorded deliberately.

\term{Oracle} An independent, direct (slow, uncompressed) evaluation of a quantity
  used as the reference in a test --- e.g.\ the direct Wannier-sum vertex against the
  SVD-tabulated one. Lesson from P8: an oracle that shares a bug with the code under
  test validates nothing; oracles must be derived independently.

\term{Round-trip} Reconstructing upstream ground truth from the decomposed parts:
  $g^{\mathcal S}+\langle u|u\rangle g^{\mathcal L}$ against the coarse \texttt{.epb}
  coupling (machine precision at all $215$ coarse $\qv$), and short-range $+$
  long-range dynamical matrix against \texttt{dynq}. The strongest data-preparation
  check in the repo.

\term{Executable spec} A deliberately flat, fixed-capacity CPU re-implementation of
  the four sampler updates, simple enough to audit line-by-line, used as the
  reference the GPU kernel must match bit-for-bit (stages (i)--(ii) of the GPU
  validation sequence).

\term{Parity} Agreement between two implementations at a stated strictness:
  \emph{bit parity} (identical reports to the last bit --- GPU kernel on the CubeCL
  CPU runtime vs the flat driver; multi-GPU merge vs single-device run) or
  \emph{accept-count parity} (identical acceptance counts on real hardware, where
  lane scheduling forbids bitwise transcript equality). \texttt{ironmc gpu parity}
  runs the hardware check.

\term{Scout-sized vs production} Two sizes of the same measurement: a scout run
  (minutes, larger error bars) to steer, then the production run (node-hours, quoted
  statistics) to decide. Gates default to scout size and take an environment override
  for certification.

\term{Anchor} An external reference number a validation stage lands on (Mishchenko
  $\alpha=1,2$; FEP-DMC $-0.2479$~eV at $20^{3}$).
\end{description}

\subsection{Codebase}

\begin{description}
\term{\texttt{ironmc}} This Rust workspace: DiagMC engines and batteries for polaron
  problems.

\term{\texttt{rmc-core}} The engine crate: updates, Metropolis kernels, update sets,
  runner, seeding.

\term{\texttt{rmc-diagmc}} The state-generic DiagMC engine (runner, updates, decay
  estimator) behind the \texttt{Model} trait.

\term{\texttt{Model} trait} The abstraction a material implements: electron
  dispersion $\varepsilon(\kv)$, phonon dispersion $\omega_{\qv\nu}$, complex vertex
  $g(\kv,\qv,\nu)$, and a phonon-momentum proposal. Two implementations: analytic
  Fr\"ohlich and the ab-initio LiF model.

\term{\texttt{rmc-frohlich}} The standalone Fr\"ohlich polaron engine and performance
  fixture.

\term{\texttt{rmc-ab-initio}} The ab-initio consumer app: EPW readers, Wannier
  reconstruction, polar terms, SVD tabulation, table building, and the production
  run/merge CLI.

\term{\texttt{rmc-minimal}} The hot-path benchmark fixture crate.

\term{Config (\texttt{configs/*.json})} A run definition: sampler knobs
  (\texttt{mu}, \texttt{max\_tau}, \texttt{chains}, \texttt{max\_steps},
  \texttt{warmup\_steps}, \texttt{n\_batches}, \texttt{seed},
  \texttt{fixed\_tau}, \dots). Naming: \texttt{stage1-*}/\texttt{stage2-*} for the
  validation stages; \texttt{ladder-lif*} for grid-ladder production
  (\texttt{-node} $=$ full-node chain count, \texttt{-tau120} $=$ guard-A variant,
  \texttt{-gpu} $=$ GPU population, \texttt{-smoke} $=$ smoke-sized).

\term{Report} The JSON a run emits (\texttt{e0\_thermo}, \texttt{sign},
  \texttt{mean\_order}, \texttt{steps\_done}, \dots), consumed as structured data by
  the \texttt{nu} tooling.

\term{\texttt{RankDump} / \texttt{merge}} Per-SLURM-rank raw measurement dump written
  by multi-GPU jobs (default under \texttt{results-gpu/<jobid>}), merged offline into
  one report; the merge is bit-identical to the equivalent single-device run by the
  global-chain-id construction.
\end{description}

\subsection{Reproducibility and RNG}

\begin{description}
\term{Deterministic chain streams} The load-bearing invariant: each chain's RNG
  stream is derived from (master seed, \emph{stable} chain id), never from scheduling
  order --- so thread/rayon scheduling can never change which stream a chain gets.
  CPU: xoshiro256++ streams. GPU: counter-based Philox4x32-10 keyed by
  (seed, chain, step), which lets all $32$ lanes of a warp compute identical draws
  with no communication.

\term{Global chain id} Chain ids are global across ranks (offset $=$
  rank $\times$ chains-per-rank), so a multi-rank population is a pure repacking of a
  single larger run and merged statistics are a function of (seed, total chains)
  only.

\term{f64 policy} All physics arithmetic is f64 (weights, $\tau$, energies);
  f32 appears only as a \emph{storage} format for the big factor tables. Momenta are
  exact integer grid indices.
\end{description}

\subsection{GPU port}

\begin{description}
\term{CubeCL} The Rust GPU compute framework used for the CUDA backend; its CPU
  runtime executes the same kernel on the host, enabling the bit-parity stage.

\term{Warp-per-chain / redundant lanes} The kernel's parallel mapping: one chain per
  $32$-lane warp, every lane redundantly executing the chain's update logic
  (warp-uniform control flow for free), lanes cooperating only inside vertex
  evaluations.

\term{Lockstep updates} All lanes of a chain's warp (and, under group-uniform
  selection, all chains of a group) attempt the same update type each step. Mandatory
  on device: barrier arrivals pair by dynamic count, so divergent lanes either
  deadlock or silently corrupt the shared scratch memory --- and the barrier count per
  step is state-dependent, so uniformity must hold over the whole dynamic trace.
  Statistically free: the update-type schedule is state-independent, so each chain's
  marginal law is the standard random-scan chain (selection probabilities cancel in
  the acceptance ratio as usual); a shared schedule only correlates chains weakly ---
  variance, never bias. Corollary: impossible proposals count as rejections, never
  redrawn (a state-dependent redraw would bias the run).

\term{Cooperative vertex dot} The rank-$r$ contraction strided across the $32$ lanes
  with a shared-memory reduce; the emit/absorb pair
  $g(\kv_{e},\qv,\nu)\,g^{*}(\kv_{a},\qv,\nu)$ fuses both dots over one shared
  $v_{g}(\qv,\nu)$ strip, halving table traffic.

\term{Bandwidth-bound} The regime the design targets: a vertex evaluation streams
  $\approx16r$ bytes, so HBM bandwidth ($\approx1.24$~TB/s sustained on an A100,
  $\approx80\%$ of peak) --- not FLOPs --- sets throughput; this is where a GPU beats a
  CPU node by an order of magnitude, and why the margin grows with grid size once the
  table no longer fits CPU L3.

\term{Planar layout} Complex tables stored as separate real/imaginary f32 planes for
  coalesced loads.

\term{\texttt{ironmc gpu} commands} \texttt{run} (production GPU run),
  \texttt{scale} (chain-count throughput sweep), \texttt{parity} (hardware parity
  test), \texttt{vbench} (vertex-evaluation microbenchmark), \texttt{info}.
\end{description}

\subsection{Cluster operations}

\begin{description}
\term{VSC-5} The Vienna Scientific Cluster the campaigns run on: $128$-core Zen3 CPU
  nodes ($512$~GB--$2$~TB) and $2\times$A100 GPU nodes, driven via SLURM.

\term{Partition / QoS / devel} SLURM node classes (\texttt{zen3\_0512},
  \texttt{zen3\_0512\_a100x2}, \dots) with default and \emph{devel} QoS --- the devel
  queue gives fast turnaround at a fixed $10$-minute walltime cap; smoke- and
  scout-sized runs are shaped to fit it. The partition table (including GPU counts)
  lives in \texttt{nu/ironmc/slurm.nu}.

\term{\texttt{ironmc slurm submit}} Wraps a target command (e.g.\
  \texttt{ladder run 50}) in a generated per-job SLURM script with its own job
  directory and results under \texttt{./results/}.

\term{Deadline truncation} Sizing \texttt{max\_steps} beyond the walltime and letting
  the SLURM end-time stop the run; statistics remain valid because batches pool
  across chains and throughput is quoted from \texttt{steps\_done}.

\term{\texttt{--node} config} The full-node CPU variant of a ladder run: chain count
  auto-sized to the CPU allocation, $10^{7}$ steps per chain.
\end{description}

\subsection{Tooling and process}

\begin{description}
\term{\texttt{rtk}} Token-optimizing CLI proxy; shell commands in agent sessions are
  prefixed with \texttt{rtk} (e.g.\ \texttt{rtk git status}).

\term{\texttt{nu/ironmc}} The Nushell operations module (\texttt{ladder},
  \texttt{tables}, \texttt{gates}, \texttt{gpu}, \texttt{slurm}, \texttt{qe},
  \texttt{bench}, \texttt{test} submodules); every command returns structured records
  meant to be piped, not parsed.

\term{\texttt{journal.md}} The append-only lab notebook: dated entries recording each
  implementation chunk, calibration, gate result, and lesson. Methodology graduates
  from the journal into the LaTeX note.

\term{Plans (\texttt{plans/*.md})} Design documents for work campaigns
  (\texttt{plan-large-grids.md}, \texttt{gpu-abinitio.md}, \dots). Steps inside a
  plan are labeled \texttt{A0, A1, \dots} (phases), so ``A0'' or ``D7'' always refers
  to a step \emph{of a specific plan}; the journal names the plan when ambiguous.

\term{G/N/D/P labels} The methodology note's own numbering: \textbf{G}oals,
  \textbf{N}on-goals, \textbf{D}eferred items, \textbf{P}itfalls (P1--P9). Retired
  non-goals keep their labels for reference.

\term{Phase A / Phase B} The two halves of the original ab-initio plan: A $=$
  data preparation (QE/EPW runs, readers, reconstruction), B $=$ the DiagMC consumer
  (sampling the prepared vertex).

\term{Profiles} Cargo build profiles: \texttt{release} is the default for everything
  (debug is off limits without a stated reason); \texttt{release-tuned} exists for
  tuned performance builds.

\term{\texttt{bench-compare}} The benchmark harness alias comparing steps/s across
  revisions with repetitions; benchmarks need an isolated core and are run by the
  user, not by agents.
\end{description}

\subsection{Overloaded terms --- read by context}
\label{sec:overloaded}

\begin{description}
\term{rank} (1) SVD rank budget of the compressed vertex ($800$, $1600$) --- the sense
  in \texttt{--rank}, \texttt{gate-dmc}, $\Delta_{\mathrm{rank}}$, table filenames.
  (2) SLURM/MPI rank --- one process, one GPU, in multi-device jobs (\texttt{RankDump},
  ``two-rank job''). Table names use sense (1); results directories use sense (2).

\term{rung (retired)} Historical term, removed from the working tree on 2026-09-11
  but surviving in git history and commit messages. It entered via the
  \texttt{plan-rmc-ab-initio.md} ``validation ladder'' metaphor and ended up carrying
  two senses at once --- a validation \emph{stage} (``rung 1'' $=$ Fr\"ohlich engine
  check, ``rung 2'' $=$ real-tensor check, ``rung-2 checkpoint'') and a grid-ladder
  \emph{point} (``rung-1 run'' in throughput contexts $=$ the $20^{3}$ production
  run) --- which is why it was retired. Reading old history: validation contexts
  $\to$ \emph{stage}, grid contexts $\to$ \emph{point}.

\term{ladder} Means the grid ladder only (the $1/n$ extrapolation campaign; the
  \texttt{ironmc ladder} command family). Pre-rename documents also said
  ``validation ladder'' for what is now \emph{validation stages}. For external
  writing, prefer ``grid-convergence sequence'': in many-body theory ``ladder
  diagrams'' and the ladder approximation are established, unrelated terms.

\term{vertex} (1) The e--ph matrix element $g(\kv,\qv,\nu)$ (``vertex evaluation'',
  ``vertex cache''). (2) A diagrammatic interaction event on the worldline (``vertex
  count $N_{v}$'').

\term{order} (1) Diagram order (number of phonon arcs). (2) Perturbative order
  ($E^{(2)}$ is second order in $g$). (3) Ordering of $\tau$ values in the diagram
  array (``$\tau$-sorted'').

\term{table} (1) The \texttt{.rmcvtab} discretized-Hamiltonian artifact. (2) The
  in-memory SVD factor tables. (3) Nushell's tabular output.
\end{description}

\end{document}
